\papersection{Sensibility}{sec:sense}

Markets rarely react to the same news on the same clock. The shape of the repricing repeats, but
the timing may differ by a few minutes from one release to the next. A distance metric that cannot
see through that slippage will call two identical reactions different. Dynamic time warping (DTW)
measures the similarity between two time series while letting their time axes stretch and
compress, so two paths are close if they trace the same shape, even when one runs faster or lags
the other \citep{Sakoe1978, Berndt1994}. This feature makes DTW well suited to price data:
Euclidean distance compares minute by minute, and correlation is just as rigid in time. We use
DTW as an analogue forecaster \citep{Lorenz1969}. We rank historical price paths by their
similarity to a target window and read off the returns that followed. We posit that similar price
paths have similar short-term forward returns. Figure~\ref{fig:dtw} illustrates the alignment.

\begin{figure}[!tb]
  \centering
  \includegraphics[width=0.6\linewidth]{dtw_illustration}
  \caption{How DTW compares two price paths (stylized example). Panel (a): the lines connect the
  points that DTW aligns; the past path reaches the same low about two minutes later. Panel (b):
  the local cost inside the band $|i-j|\le 5$ and the optimal warping path (orange); the dashed
  diagonal is the minute-by-minute (Euclidean) alignment.}
  \label{fig:dtw}
\end{figure}

Applied naively, DTW runs into two problems. The first is cost: comparing a target against every
window in years of minute-level data is slow. The second is that, with a large enough library,
some historical paths will match the target by pure luck, and a lucky match says nothing about why
the path formed. We address both issues by restricting the library to windows that follow
scheduled macroeconomic releases such as CPI and the Jobs Report, events macro traders reliably
react to. The library shrinks, and every match shares a common root: one public shock at a known
time, with participants repricing rate expectations together. Our refined hypothesis is that
similar price paths following scheduled macro releases have similar short-term forward returns.
For example, take the price path following a CPI release. Compared against past paths following
other releases, such as the Jobs Report, retail sales or PCE, most will strongly differ, yet a
few will trace nearly the same shape. Those few are far more likely to reflect the same
repricing than sheer luck.
